\documentclass[a4paper,11pt]{article}
\usepackage{fancyhdr,graphics,graphicx,xy}
\usepackage{blkarray}
\usepackage{ifthen}
\usepackage{amssymb,amsmath,mathtools,color}
\pagestyle{fancy}
\usepackage{xparse}
\usepackage{nicematrix}


\textwidth = 6.5 in
\textheight = 9 in
\oddsidemargin = 0.0 in
\evensidemargin = 0.0 in
\topmargin = 0.0 in
\headheight = 0.0 in
\headsep = 0.0 in
\parskip = 0.2in
\parindent = 0.0in


%%%%%%%%%%%%%%%%%%%%%%%%%%%%%%%%%%%%%%%%%%
%%% General math macros
%%%%%%%%%%%%%%%%%%%%%%%%%%%%%%%%%%%%%%%%%%


\newcommand{\asstnumber}{2}
\newcommand{\assigntop}{
  \fancyhead[L]{}
  \setlength{\headheight}{37.40004pt}
  \fancyhead[C]{\Large Math 576 -- Fall 2023\\ Assignment \asstnumber}
  \fancyhead[R]{ }
}

\newtheorem{theorem}{Theorem}
\newtheorem{corollary}[theorem]{Corollary}
\newtheorem{definition}{Definition}

\begin{document}
\assigntop

\section{question 1}
\subsection{(a)}
We have events, $U_1$ and $U_2$ corresponding to the two bins with probability $\frac{1}{2}$ for both, and event $B$ for choosing a blue marbel, and event $Y$ for choosing a yellow marble. We choose a bin at random and then we draw a marble from the bin.
We want to know
$$\mathbb{P}(U_1|Y) = \frac{\mathbb{P}(Y | U_1) \cdot \mathbb{P}(U_1)}{\mathbb{P}(Y)}$$
knowing that 
$$
\mathbb{P}(Y | U_1) = \frac{2}{5}
$$
and
$$
\mathbb{P}(U_1) = \frac{1}{2}
$$
and
\begin{align*}
	\mathbb{P}(Y) &= \mathbb{P}(U_1 Y) + \mathbb{P}(U_2 Y)\\
		      &= \mathbb{P}(U_1) \mathbb{P}(Y | U_1) + \mathbb{P}(U_2) \mathbb{P}(Y | U_2) \\
		      &= \frac{1}{2} \frac{2}{5} + \frac{1}{2} \frac{4}{7}\\
		      &= \frac{1}{5} + \frac{2}{7} \\
		      &= \frac{17}{35}
\end{align*}
thus 
$$\mathbb{P}(U_1|Y) =\frac{ \frac{2}{5} \cdot \frac{1}{2}}{\frac{17}{35}} =\frac{7}{17} $$
\subsection{(c)}
We want (subscript $2$ denotes the second draw)
\begin{align*}
	\mathbb{P}(U_1|Y_1 Y_2) &= \frac{\mathbb{P}(Y_1 Y_2 | U_1) \mathbb{P}(U_1)}{\mathbb{P}(Y_1 Y_2 | U_1) \mathbb{P}(U_1) + \mathbb{P}(Y_1 Y_2 | U_2) \mathbb{P}(U_2)}\\
				&= \frac{2/5 \cdot 2/5 \cdot 1/2}{2/5 \cdot 2/5 \cdot 1/2 + 4/7 \cdot 4/7 \cdot 1/2} = \frac{49}{149}
\end{align*}
\clearpage
\section{2}
let $i$ be the $i$-th candidate interviewed, we basically want to compute
\begin{align}
P(k) &= \sum_{i=1}^n \mathbb{P}( i \text { best } \cap i \text{ selected } ) \\
       &= \sum_{i=1}^n \mathbb{P}( i \text{ selected } | i \text{ best } )c \cdot \mathbb{P}(i \text{ best } )\\
       &=\sum_{i=1}^n \mathbb{P}( i \text{ selected } | i \text{ best } )c \cdot \frac{1}{n}\\
       &= \sum_{i=k+1}^n \mathbb{P}( \text{ best of $i-1$ applicants amongst $k$ first } ) \cdot \frac{1}{n} \\
       &= \sum_{i=k+1}^n  \frac{k}{i-1} \cdot \frac{1}{n}\\
       &= \frac{k}{n}\sum_{i=k}^n \frac{1}{i}
\end{align}
Using Mathematica, the peak of the plot of the probability above against $k$ for some fixed $n$ peaks at around $37$ percent, which is close to $\frac{1}{e}$, with $k$ having the same ratio with $n$, thus we should pick $k^* = \frac{n}{e}$. 
\clearpage
\section{3}

Flip the coin twice each time, and map the outcomes $HT \to 0$ and $TH \to 1$ and discard the outcomes $HH$ and $TT$, thus the sample space is $\Omega = \{0,1\}$, where $\mathbb{P}(0) = p(1-p) = \mathbb{P}(1)$ so since the latter outcomes are equally as likely, it simulates a fair coin, since all flips are independent.
\clearpage
\section{4}
\subsection{(a)}
Assuming $X$ is discrete then
\begin{align*}
\mathbb{E}(u(X) + v(X)) &= \sum_{x \in S_x} (u(x) + v(x)) \mathbb{P}(X = x)\\
			  &= \sum_{x \in S_x} u(x) \mathbb{P}(X = x) + \sum_{x \in S_x} v(x)\mathbb{P}(X =x)\\
			  &= \mathbb{E}(u(X)) + \mathbb{E}(v(X))
\end{align*}
\subsection{(b)}
Assuming $X$ has a probability density function then 
$$
\mathbb{E}(u(X)) = \int_{-\infty}^{+\infty} u(x)f(x)
$$
where $f(x)$ is the probability density function of $X$.
then
\begin{align*}
	\mathbb{E}(u(X)+v(Y)) &= \int_{-\infty}^{+\infty} (u(x) + v(x))f(x) \\
			      &= \int_{-\infty}^{+\infty} u(x) f(x) +  \int_{-\infty}^{+\infty} v(x) f(x)\\
			      &= \mathbb{E}(u(X)) + \mathbb{E}(v(X))
\end{align*}
\clearpage
\section{5}
$$X = \text{number of people who receive the correct hat}$$
We need to find 
$$\mathbb{P}(X = \ell)$$ 
where $\ell \in \{1,2,\ldots,n\}$. Recall the probability that $l$ people out of $n$ gets their hat and no one else does is 
$$
f(l)=\sum_{k=1}^{n-l} \frac{(-1)^k}{k!}(l!)^{-1}
$$
which is the probability density function of $X$. Thus finding the expected value
\begin{align*}
	\mathbb{E}(X) &= \sum_{i = 1}^{n} i f(i)\\
		      &= \sum_{i=1}^n i \sum_{k=1}^{n-i} \frac{(-1)^k}{k!}(i!)^{-1}
\end{align*}
This is the expected number of fixed points in a random permutation. However this formula is pretty complicated to decifer, we can bypass this problem by actually setting
$$Y_i = \text{ i is fixed}$$ where $Y_i = 1$ if it's $i$ is fixed, and $0$ otherwise.
thus the probability mass function of $X_i$ 
$$f_i(1) = \mathbb{P}(Y_i = 1) = \frac{(n-1)!}{n!} = \frac{1}{n}$$
and 
$$E(Y_i) = \sum_{x \in S_{Y_i}} x \cdot \frac{1}{n} = \frac{1}{n}$$
since $S_{Y_i} = \{1\}$

$$X(\sigma) = Y_1(\sigma) + Y_2(\sigma) + \ldots + Y_n(\sigma)$$
For any permutation $\sigma$, which is an event in our probability space,
then the expected value should be
$$
\mathbb{E}(X) = \mathbb{E}(Y_1) + \ldots + \mathbb{E}(Y_n) = \frac{1}{n} + \ldots + \frac{1}{n} =  1
$$

\end{document}
